% Compile with: latexmk -xelatex -outdir=build main.tex
\documentclass[aspectratio=169,10pt]{beamer}
\usetheme{nvidialight}

\title{Making Sparser MoE Practical}
\subtitle{Router and expert-parallel optimizations}
\author{Your name \enspace|\enspace Your team}
\date{Internship final presentation \enspace|\enspace July 2026}
\NvidiaSetFooter{Internal presentation}

\begin{document}

{\setbeamertemplate{footline}{}\begin{frame}[plain]
  \titlepage
\end{frame}}

\NvidiaTransitionSlide{The performance problem changed}

\begin{frame}{Target workload: wide routing, small expert compute}{Why routing became the critical path}
  \NvidiaSource{Source: investigation configuration, EP72 training setup}
  \begin{columns}[T,onlytextwidth]
    \begin{column}{.53\textwidth}
      \begin{itemize}
        \item 2,304 experts and top-$k=36$ widen routing metadata.
        \item 72 expert-parallel ranks turn routing into distributed data movement.
        \item 512-wide experts make grouped GEMMs latency-sensitive.
      \end{itemize}
    \end{column}
    \begin{column}{.42\textwidth}
      \renewcommand{\arraystretch}{1.3}
      \begin{tabularx}{\linewidth}{@{}>{\NvidiaSansMedium}l X@{}}
        \toprule
        Dimension & Value \\
        \midrule
        Total experts & 2,304 \\
        Router top-$k$ & 36 \\
        Expert parallelism & 72 \\
        Local experts/rank & 32 \\
        Latent hidden size & 512 \\
        \bottomrule
      \end{tabularx}
    \end{column}
  \end{columns}
\end{frame}

\begin{frame}{One representation must survive every ownership boundary}{Avoid expanding routes at API boundaries}
  \NvidiaSource{Source: final-presentation content note}
  {\centering
    \resizebox{\textwidth}{!}{\begin{tikzpicture}[node distance=0.14in, every node/.style={font=\scriptsize,align=center,text width=1.18in,minimum height=.64in}]
      \node[draw=NvidiaRuleGray,fill=NvidiaLightGray] (router) {TE fused router};
      \node[draw=NvidiaRuleGray,fill=NvidiaLightGray,right=of router] (mcore) {Megatron Core};
      \node[draw=NvidiaRuleGray,fill=NvidiaLightGray,right=of mcore] (hybrid) {HybridEP route exchange};
      \node[draw=NvidiaRuleGray,fill=NvidiaLightGray,right=of hybrid] (gemm) {Grouped GEMMs};
      \draw[-{Latex[length=2mm]},line width=.8pt,NvidiaGreen] (router) -- node[above,font=\scriptsize,text=NvidiaGray]{selected IDs / weights} (mcore);
      \draw[-{Latex[length=2mm]},line width=.8pt,NvidiaGreen] (mcore) -- node[above,font=\scriptsize,text=NvidiaGray]{compact $[T,K]$ map} (hybrid);
      \draw[-{Latex[length=2mm]},line width=.8pt,NvidiaGreen] (hybrid) -- node[above,font=\scriptsize,text=NvidiaGray]{expert-contiguous tokens} (gemm);
    \end{tikzpicture}}\par}
  \vspace{0.03in}
  \begin{NvidiaEvidence}{Design rule}
    A compact route representation only pays off if it crosses the TE, Megatron Core, and HybridEP API boundaries without expanding back to $[T,E]$.
  \end{NvidiaEvidence}
\end{frame}

\begin{frame}{Router and HybridEP optimizations improved the matched stack}{Each result uses its original benchmark scope}
  \NvidiaSource{Source: benchmark notes; do not compare across benchmark suites}
  \noindent\NvidiaMetric{10$\times$+}{large-$E$, large-$K$ router kernel speedup}{Isolated router result}\hfill
  \NvidiaMetric{53.1\%}{matched end-to-end gain}{June benchmark set}\hfill
  \NvidiaMetric{106.6\%}{later MoE-module-only gain}{Different test scope}
  \vspace{0.03in}
  \begin{NvidiaEvidence}{Interpretation}
    Do not add isolated kernel speedups to predict training throughput: kernels overlap, metadata is reused differently by phase, and the critical path moves after every fix.
  \end{NvidiaEvidence}
\end{frame}

\NvidiaTextImageFrame{Keep evidence and takeaway together}{Use a focused visual beside a concise interpretation}{assets/nvidia-cover-background.jpg}{Source: NVIDIA light presentation template}{%
  {\NvidiaSansMedium\color{NvidiaGreen}The visual answers what happened.}\par
  The copy explains why the visual matters to the audience.
  \vspace{.08in}
  \begin{itemize}
    \item Use one clear focal point.
    \item Keep the copy to one claim and two implications.
  \end{itemize}
}

\NvidiaImageTextFrame{Image-left layouts add context}{The image leads; the copy provides meaning}{assets/nvidia-transition-background.jpg}{Source: NVIDIA light presentation template}{%
  {\NvidiaSansMedium\color{NvidiaGreen}Use this for context-first storytelling.}\par
  The visual establishes the setting before the speaker explains the consequence.
  \vspace{.08in}
  \begin{itemize}
    \item Prefer one image, not a collage.
    \item Use a 1:1 or 4:5 crop for more visual weight.
  \end{itemize}
}

\NvidiaWideImageFrame{Wide images set the context}{Use the 16:9 one-up pattern for a single visual}{assets/nvidia-cover-background.jpg}{Source: NVIDIA light presentation template}{Use a high-resolution 16:9 image so the visual remains legible at presentation distance.}

\NvidiaTransitionSlide{Thank you}

\end{document}
